%
\begin{frame}[t]%
\sharedframetitle{Components needed for Metaheuristic Optimization}%
%
\begin{itemize}%
%
\item We have a space~\searchSpace\ in which we search for potential solutions~\sespel.%
%
\item<2-> \searchSpace~is often very/exponentially large in the input size.%
%
\item<3-> We repeatedly sample elements~\sespel\ from this space~\searchSpace, so we need operators to do that\only<-3>{.}\only<4->{:%
\begin{itemize}%
\item a \emph{nullary} operator usually returns a randomly chosen point from~$\searchSpace$.\uncover<5->{ %
This is often used to get starting points for the search.}%
%
\item<6-> a \emph{unary} operator receives an existing point as input and creates a slightly modified copy.\uncover<7->{ %
This is often used to explore the neighborhood of good solutions.}%
%
\item<8-> a \emph{binary} operator receives a two existing points as input and creates a new point with elements from both of these \inQuotes{parents}.\uncover<9->{ %
This is often used to combine traits from good solutions.}%
%
\end{itemize}%
}%
%
\item<10-> We evaluate the quality of solutions using the objective function~$\objFun:\searchSpace\mapsto\realNumbers$, which usually is subject to minimization~(smaller values are better).%
%
\item<11-> These are the usual main ingredients of \gls{metaheuristic} optimization.%
%
\end{itemize}%
%
\end{frame}%
%
\global\let\sharedSlideTitle\relax%
%
